\documentclass[11pt]{article}
\usepackage[margin=1in]{geometry}
\usepackage{amsmath,amssymb,amsthm,hyperref}
\newtheorem{proposition}{Proposition}
\newtheorem{lemma}{Lemma}
\title{Hillclimb 002: the actual exponent target and an executable C38 graph}
\author{Research audit; no new ECDLP complexity bound claimed}
\date{31 October 2026}
\begin{document}
\maketitle

\paragraph{Status and sources.}
The preceding iteration's coefficient $27.13327$ remains conditional.
The inherited uniform contraction theorem, arithmetic proof, and complete
ECDLP transfer have not been independently certified. This iteration
establishes a barrier for a specified enumeration strategy, supplies
independent finite Smith checks, and implements the previously missing
explicit C38 graph interface. The source hashes are recorded in the run
records and \texttt{spectrum\_results.json}. Original archives are unchanged.

\section{What reaching the stated target actually requires}
The goal is classical expected time $2^{\widetilde O(\sqrt b)}$ for
general prime-field ECDLP in the stated supplied-order and cofactor class.
The complete paid relation reduction currently has the form
\[
 T\le\operatorname{poly}(b,m)(BC_m+mB^2),\qquad
 \log_2B=b/m+O(\log b).
\]
The arithmetic network uses $S=\widetilde O(mb)$ edges. Replacing its
contraction algorithm changes $C_m$; it does not automatically remove
the factor-base construction or the $B^2$ term.

\begin{proposition}[Polynomial-exponent transfer within this reduction]
Suppose a contractor with $\log_2 C_m=\widetilde O(S^\theta)$,
$0<\theta\le1$, is valid uniformly on the arithmetic graphs at the
chosen arities, and suppose the entire relation-reduction proof also
extends to those arities. Setting $m=b^\gamma$ gives the certified
power of $b$ in the complete exponent
\[
 \max\{1-\gamma,\ \theta(1+\gamma)\}.
\]
Its minimum for $0\le\gamma\le1$ is
\[
 \frac{2\theta}{1+\theta},\qquad
 \gamma=\frac{1-\theta}{1+\theta}.
\]
Thus a square-root-in-$S$ contractor would yield power $2/3$ here,
whereas power $1/2$ via this particular accounting requires
$\theta\le1/3$, with $m=b^{1/2}$ at equality, up to logarithmic factors.
\end{proposition}
\begin{proof}
The factor-base exponents are $b^{1-\gamma}$; the relation exponent is
$b^{\theta(1+\gamma)}$, up to logarithmic factors. Multiplication adds
exponents but does not change the larger polynomial power. The first
displayed function decreases and the second increases with $\gamma$.
Solving their equality proves the formula.
\end{proof}

This optimizes a proposed upper-bound formula; it is not a lower bound
on ECDLP or on every implementation of the reduction. In particular,
an algorithm that avoids explicit factor-base enumeration may have a
different balance. The uniform arity hypotheses in the proposition
are essential. The existing composition was stated for polylogarithmic
$m$. At $m=b^{1/2}$ its leading degree ledger becomes roughly
$\Delta=\widetilde O(m\sqrt b)=\widetilde O(b)$ and
$S=\widetilde O(b^{3/2})$, outside the low-degree condition
$\Delta\le S^{1/2}\operatorname{polylog}S$. A proposed new contractor
and the coverage proof must be audited in this different regime.

For a more specialized hypothetical relation exponent
$\widetilde O(m^a b^c)$ with $a\ge0$, $0\le c\le1$, the same
balance gives power $(a+c)/(1+a)$ at
$\gamma=(1-c)/(1+a)$. Reaching power at most $1/2$ requires
$a+2c\le1$ within this model. A route may instead change the reduction.

\section{A constructive barrier for unrestricted mode enumeration}
The inherited crossing block has normalized Smith exponents
\[
 2q+2v_2(q!)=4q-2s_2(q),\quad 0\le q<D/2,
\]
with repeated residue families. Restoring the common deformation factor
adds a shift $\delta$. In the uniform schedule one can take
$\delta=\log_2D+1$, or the weaker shift one for upper bounds.
The precision cost is $P=(6+o(1))MD/u$, and the number of families
is of order $M/D$. These are the three linked quantities that produce
the square-root partition estimate.

\begin{lemma}[Many admissible tuples even with a degree cap]
Consider $Q=M/D$ families, with integer $Q,D$, $u\ge4$, and weights
\[
 w_q=\delta+4q-2s_2(q),\qquad 0\le q<D/2.
\]
Put $r=\lceil\sqrt u\rceil$, $k=\lfloor D/r\rfloor$.
Assume $D\ge2r$, $D\ge\delta\sqrt u$, and $P>3MD/u$.
There are at least $2^{M/(4\sqrt u)}$ tuples with weight less
than $P$ and total degree at most $M$.
\end{lemma}
\begin{proof}
Use only ranks $0,\ldots,k-1$ in every family, and allow each of
their exponents independently to be zero or one. There are $2^{Qk}$
choices. Their maximum degree is $Qk\le M$, and their maximum weight is
\[
 Q\sum_{q<k}w_q
 \le Q(\delta k+2k(k-1))
 \le \frac{M\delta}{\sqrt u}+\frac{2MD}{u}
 \le\frac{3MD}{u}<P.
\]
Also $k\ge D/(2r)$ and $r\le2\sqrt u$, so
$Qk\ge M/(4\sqrt u)$. Finally $k\le D/2$ since $u\ge4$.
\end{proof}

The schedule's polynomially large $D$ satisfies these guards
asymptotically. Constant fractions of the family universe give the
same power with changed constants. Hence explicit enumeration of the
unrestricted global or fixed-fraction projected candidate universe
cannot improve its $M/\sqrt u$ power. The padding/label schedule has
$u=O(\log S)$, retaining the existing logarithmic scale.

\emph{Scope of this barrier:} the tuples just constructed are feasible
indices, not proved nonzero tensor coefficients. Empty support, arithmetic
identities, cancellations, batched evaluation or symbolic summation may
eliminate or avoid them. This is not a lower bound on actual C38 support,
tensor contraction, or ECDLP. It specifically rejects claims that a new
enumeration order or constant-factor matrix improvement changes this
power while still visiting that universe.

\paragraph{Finite certificates.}
The checker independently constructs the rational Cauchy matrices
$2^i/(T+x-y)$ on equal-bit slots. It obtains their Smith exponents by
minimum-valuation pivots and integral elementary operations over
$\mathbb Z_{(2)}$, without importing the source's explicit mode formula.
All 60 scalar cases through $L=32$ agree. These checks do not establish
the extension-ring or nilpotent assertions at arbitrary sizes.
Six explicit large index-universe certificates use
$D=2^{2u}$, $M=2^{10u}$ for $u=4,8,16,32,64,128$.
The exact sum of digit counts evaluates the weight of all selected modes
without enumerating exponentially many indices. The count is certified
by the independent binary choices above, not by a floating-point fit.

\section{An additional consequence that the foundation must support}
The inherited theorem is stated for all succinct Boolean networks with
$\Delta\le S^{1/2}\operatorname{polylog}S$ and polylogarithmic parallel
multiplicity, not just for elliptic-curve predicates.

Here is a direct consequence of that claimed scope. Given a CNF formula
with clauses of width at most three and $N$ literal occurrences, create
one vertex per occurrence and one per clause. Connect each occurrence
to its clause. For every variable, connect its occurrence vertices by
a path. Each occurrence tensor requires all its incident bits to be
equal; its degree is at most three. Each clause tensor evaluates the
OR of its incident literals with their specified signs; its degree is
at most three. There are at most $2N$ edges, no loops, and no parallel
edges. Every satisfying assignment to the variables occurring in the
formula corresponds to exactly one accepting edge assignment. Variables
with no occurrences contribute a known power of two.

Therefore the inherited theorem, if valid as stated, also gives exact
counting for these formulas in $2^{O(N/\sqrt{\log N})}$ time. This is
a consequence derived from the stated input model, not an additional
assumption and not a proof that the theorem is false. It makes clear
that auditing the foundation is a substantive research obligation;
arithmetic fixtures alone cannot certify a theorem of this scope.

\section{The missing C38 graph interface is now executable}
\texttt{c38\_network.py} builds the general-prime graph from public
parameters and constants only. It creates one descriptor for every first,
middle and closure tensor, plus every binary equality relay. Fourier
words are jointly recoded in the root-fixed simplex alphabet, with
separate edge identifiers for every symbol. Each side-state relay has
two distinct edges. The source and sink vertices independently read
their own copy, and the relay enforces equality. Initialized and terminal
range states are constants, not free edges.

Tensor entry programs receive only incident symbols. A checked local
reader raises \texttt{NonincidentRead} on an undeclared read; this error
is not treated as a rejected assignment. Thus a hidden dependency would
fail the diagnostic rather than silently impose an invalid predicate.
Predicates decode their own words, reject noncanonical Fourier values
and nonbinary auxiliary symbols, and evaluate the inherited arithmetic.
They never capture a planted witness or receive a complete assignment.

The explicit ledger agrees with the inherited ledger on the tested
inputs. The compiler implements the general-prime branch only; it does
not silently claim the separate sparse-modulus fanout construction.
Witness serialization is a separate testing helper. Pinning creates
local restricted descriptors and removes only fixed edges and evaluated
scalar-one tensors. An inconsistent fixed tensor leaves an explicit
scalar-zero factor. Conditional counting recovers each remaining symbol
and then decodes and verifies the resulting curve tuple.

\paragraph{What the executed contractions establish.}
There are 18 compiled networks, 89,892 edges across fixtures,
and 10,518 explicit equality relays. Split-relay
faults and coherent Fourier mutations are rejected. Twelve recursive
contractions and witness recoveries cover real Fourier ports, carry
states, range states, distinctness states, and both sides of a relay.
These runs explicitly pin all but one or two edges. Their integer counts
agree with full enumeration of those free symbols; all active calls in
the initial contractions also agree with direct modular summation.

This is a graph-to-contraction interface check. It does not perform an
unpinned relation search. The recursive engine still uses measured small
labels, a fixed diagnostic depth, and an explicit Strassen backend,
rather than the asymptotic packing schedule and formula catalog.

\paragraph{A separate complete small source search.}
\texttt{verify\_source\_enumeration.py} exhausts all source triples
$(x_1,x_2,\lambda_2)\in\mathbb F_{13}^3$ for each of 12 two-slot
instances whose hashes were chosen without planting a decomposition.
It compares the resulting complete graph witnesses with an independently
implemented exhaustive curve-point-pair search. Both covered and uncovered
targets occur. This checks the finite arithmetic reduction in a way that
planted witnesses alone do not. It uses enumeration, not an unpinned
network contraction, and remains a finite check.

\section{Research decisions and next falsifiable hypotheses}
\begin{enumerate}
\item Reject a power-improvement claim based solely on rearranging the
same unrestricted mode tuples. The explicit subset lemma supplies the
countercriterion. Keep only methods that avoid enumeration or prove
additional support restrictions on the actual arithmetic descriptors.
\item Test whether the C38 projected-child tables admit an exact compact
representation closed under the recursion. Acceptance requires a proved
representation-size bound, algorithms for entry access and composition,
and an exponent calculation including construction cost. Small numerical
rank or a planted/pinned fixture is insufficient. A growing family with
forced large representation size rejects the proposed representation,
not every possible compression method.
\item A new kernel/spectrum or precision profile must be evaluated jointly:
deformation degree, reconstruction multiplicities, scalar precision loss,
mode multiplicities, packet geometry and input costs must all be charged.
A larger Smith weight alone is insufficient if the precision rises too.
\item Independently audit the full theorem on succinct bounded-degree
networks, using the CNF consequence as a scope check. The small Smith and
arithmetic checks above certify only their stated finite cases.
\item Before promoting any polynomial-exponent contractor to the ECDLP
goal, apply the transfer proposition and recheck its required arity and
degree regime. Reaching $2^{\widetilde O(\sqrt S)}$ alone is insufficient
for the present factor-base route.
\end{enumerate}
The active goal remains unachieved. No new time coefficient or saving
power is claimed in this iteration.
\end{document}
